\section{Motivation and Scope}

The immediate goal is to create a controlled algebraic environment for thinking
about second-moment shocks before moving to the full small-open-economy
nonlinear BCA framework. The exercise starts from the smallest useful New
Keynesian system:
\begin{enumerate}[label=(\roman*)]
    \item an intertemporal demand equation,
    \item a price-setting equation, and
    \item a monetary policy rule.
\end{enumerate}
The only primitive uncertainty shock is stochastic volatility in productivity.
The question is whether the effects of that volatility shock can be organized as
wedges in the three equations.

This note should be read as an accounting device rather than a final structural
model. A first-order New Keynesian model is certainty equivalent, so a pure
volatility shock has no independent effect on allocations. A second-order
expansion of the underlying nonlinear equilibrium conditions produces risk
correction terms, and those terms can be collected as wedges. In a fully solved
DSGE model with stochastic volatility, third-order perturbation is usually
required for volatility shocks to enter the policy functions as independent
state variables. The second-order exercise here is therefore a diagnostic step:
it clarifies where the wedges would appear and what restrictions would be
implied by the assumption that TFP uncertainty is the only primitive source of
second-moment variation.

The intended bridge to the larger project is:
\begin{align*}
    \text{TFP volatility shock}
    &\Rightarrow
    \text{risk correction terms}
    \\
    &\Rightarrow
    \text{NK wedges}
    \\
    &\Rightarrow
    \text{BCA-style counterfactuals}.
\end{align*}

The note deliberately does not impose the later two-factor empirical structure
from the OI-SVMVAR. It focuses on a single source of uncertainty so that the
mapping from volatility to wedges is transparent.
